\documentclass[11pt]{article}
\usepackage{fontspec}
\usepackage[french]{babel}
\usepackage{amsmath}
\usepackage{amssymb}
\usepackage{amsthm}
\usepackage{mathrsfs}
\usepackage{tikz-cd}
\usepackage{enumitem}

\frenchsetup{AutoSpacePunctuation=false}

\newcommand{\printedpage}[1]{\marginpar{\footnotesize p.~#1}}
\newcommand{\msp}[1]{\marginpar{\footnotesize [#1]}}
\newcommand{\Ch}{\underline{\mathrm{Ch}}}

\begin{document}

\printedpage{9}
% head: A. Grothendieck : Dérivateurs, Ch. XII, Édition des Universités de Montpellier II et Paris VII

Ainsi, le lemme de la page 7 est prouvé, moyennant l'axiome W(4) qui assure qu'une flèche qui est à la fois $W$-lisse et $W$-propre, satisfait à la condition très forte a) des $W$-fibrations.

\textbf{NB} En fait, il suffirait, au lieu de l'axiome W(4) sous sa forme forte, de supposer qu'une $f : X \longrightarrow Y$ qui est à la fois `homotopiquement lisse' et `homotopiquement propre', i.e. les fibres des $X/x \longrightarrow Y/y$ et des $x\backslash X \longrightarrow y\backslash Y$ sont contractiles (donc $W$-asphériques pour tout $W$ satisfaisant W(1,2,3)), soit $W$-fibration forte. Énonçons quand même ce que nous avons prouvé.

\textbf{Proposition 4.} \textit{Si $W$ satisfait} W(1,2,3) (\footnote[17]{Pour W(2), il le suffit sous la forme W(2 bis).}) (nous le supposons toujours par la suite) \textit{plus} W(4) (sous une forme un peu affaiblie suffisante, cf. remarque précédente), \textit{on a ceci: Soit $f : X \longrightarrow Y$ une flèche dans} Cat$/S$\textit{, avec $X$, $Y$ des $W$-fibrés sur $S$. Si $f \in W$, alors $f_s : X_s \longrightarrow Y_s$ est dans $W$ pour tout $s \in S$.}

Je me demande si j'en sais assez pour pouvoir prouver
\[ W \subseteq W_\infty . \]
Il faut quand même éliminer des cas canularesques, comme $W = \mathrm{Fl}(\mathrm{Cat})$, donc mettre l'axiome

\begin{itemize}[leftmargin=4em, labelwidth=3em, labelsep=0.6em]
\item[W(5)] (Axiome de connexité (\footnote[18]{Voir à ce sujet le corollaire page 22.}).) Si $f : X \longrightarrow Y$ est dans $W$, $\pi_0(f) : \pi_0(X) \longrightarrow \pi_0(Y)$ \texttt{[est]} bijectif.
\end{itemize}

\msp{page 11}Je factorise alors $f$ en (\footnote[19]{À nouveau, dans ces arguments on a intérêt à travailler avec $W_\omega$.})
\[ X \xrightarrow{i} \bar{X} \xrightarrow{\bar{f}} Y, \]
avec $\bar{f}$ une $W_\infty$-fibration et une $W$-fibration (\footnote[20]{en fait, $\bar{f}$ est à la fois $W_\infty$-parfait (i.e. $W_\infty$-lisse et propre) et $W$-parfait (en fait, $W_\omega$-parfait).}), $i \in W \cap W_\infty$ - cela marche finallement, % [sic] finallement -> finalement, spelling
par la construction de VII, p. 170-171 (de \textsc{Cartan-Serre}), à partir des propiétés de % [sic] propiétés -> propriétés, missing letter
$Y \longrightarrow \Ch(Y) \longrightarrow Y \times Y$ sans plus. Comme $\bar{f} i = f \in W$ et $i \in W$, on en déduit $\bar{f} \in W$. Comme $i \in W_\infty$, on en conclut $f \in W_\infty \iff \bar{f} \in W_\infty$. Donc pour prouver que $f$ est dans $W_\infty$, on est ramené au cas où $f$ est une $W$-fibration et une $W_\infty$-fibration. Mais alors la proposition 4 nous montre que les fibres de $f$ sont $W$-asphériques (\footnote[21]{on n'a utilisé ici que cette propriété sur les fibres, comme consequence de W(4). (On utilise W(4) sous sa forme la plus faible W(4b$''$), cf. pages 13, 14).}), et il suffit de voir qu'elles sont $W_\infty$-asphériques (proposition 3 (\footnote[22]{on n'utilise la proposition 3 que dans le cas $Y = S$, pour lequel W(2 bis) suffit.}) appliquée à $W_\infty$). Ainsi on est ramené au cas où $Y = e$, i.e. à prouver que si $X$ dans Cat est $W$-asphérique, il est $W_\infty$-asphérique.
% [sic] (footnote 21) consequence -> conséquence, missing accent

Soit $x_0 \in \mathrm{Ob}\, X$, considérons comme tantôt
\[ E = \Ch(X, x_0, -) \xrightarrow{p} X, \]
% foot: Version 7 décembre 2002
\printedpage{10}
% head: A. Grothendieck : Dérivateurs, Ch. XII, Édition des Universités de Montpellier II et Paris VII
on sait que $E$ est $W \cap W_\infty$-asphérique, et que c'est une fibration pour $W$ et pour $W_\infty$. Comme $X$ est $W$-asphérique (comme $E$), $p \in W$, donc à nouveau par la proposition 4 (\footnote[23]{en fait, la condition W(4b$''$) sur $W$.}), les fibres de $p$ sont $W$-asphériques, en particulier (par W(5)) 0-connexes. On a donc prouvé le

\msp{page 12}\textbf{Lemme.} (\footnote[24]{Utilise W(1, 2 bis, 3, 4b$''$) sur $W$.}) \textit{Soit $(X, x_0)$ dans} Cat$_\bullet$\textit{. Si $X$ est $W$-asphérique, alors $\Omega(X, x_0) \overset{\text{déf}}{=} \Ch(X; x_0, x_0)$ l'est aussi} (moyennant les axiomes W(1,2,3,4) sur $W$)\textit{. Donc (par itération) les $\Omega^n(X, x_0)$ sont $W$-asphériques} ($\forall\, n \in \mathbf{N}$).

\textbf{Corollaire.} \textit{Si de plus $W$ satisfait l'axiome} W(5) \textit{de connexité, alors $X$ est $W_\infty$-asphérique.}

En effet, `on sait' que cela signifie que tous les $\Omega^n(X, x_0)$ sont 0-connexes. Or l'axiome de connexité W(5) assure que $W$-asphérique implique 0-connexe.

Mais le `on sait' est un tout petit peu léger (\footnote[25]{C'est prouvé fort simplement dans le corollaire du `theorème 1', page 190.}). Pour le justifier, il faudrait, ou bien % [sic] (footnote 25) theorème -> théorème, missing accent
que je fasse le lien avec la théorie 1/2-simpliciale classique, et prouver que les $\Omega^n$ que je construis sont `les bons', i.e. que leur nerf redonne bien (à isomorphisme près dans Hot) les $\Omega^n$ classiques (??). En particulier, que le $\pi_0(\Omega^n(X, x_0))$ est bien le $\pi_n(X, x_0)$ `de tout le monde'. Ou bien que je prouve directement que la 0-connexité des $\Omega^n(X, x_0)$ équivaut à la $W_\infty$-asphéricité de $X$, définie cohomologiquement à la \textsc{Artin-Mazur}. C'est donc un travail qu'il me faudra faire absolument. Mais ce n'est plus qu'un simple travail de routine, et je peux énoncer dès maintenant le

\textbf{Théorème 1.} (\footnote[26]{Il suffit même de W(4) sous une forme affaiblie, avec des foncteurs `fibrants' dans un sens `absolu'. Cf. le NB p. 10 - et surtout, cf. pages 13, 14.},\footnote[27]{NB Au lieu de W(5), il suffit de supposer $W \neq \mathrm{Fl}(\mathrm{Cat})$, cf. th. 2, p. 20.},\footnote[28]{Les hypothèses utiles sur $W$ sont finalement W(1), W(2 bis), W(3), W(4b$''$), et $W \neq \mathrm{Fl}(\mathrm{Cat})$.}) \textit{Soit $W \subseteq$} Fl(Cat) \textit{satisfaisant les conditions} W(1,2,3,4,5) (axiome de saturation, de localisation, de l'objet final, de fibration, de connexité)\textit{. Alors}
\[ W \subseteq W_\infty . \]

\textbf{Corollaire.} \textit{On peut caractériser $W_\infty$ comme la plus grande partie de} Fl(Cat) \textit{satisfaisant aux}
\msp{page 13}\textit{conditions} W(1) \textit{à} W(5) \textit{précédentes.}

Bien sûr, on aimerait un axiome 6 plus intéressant, qui joint aux précédents assure que $W = W_\infty$ (\footnote[29]{Cf. axiome W(6) page suivante.}). Ou est-ce que ces propriétés assurent déjà que $W = W_\infty$? Pour le voir, il faudrait peut-être que j'en sache plus sur la structure des objets $W_\infty$-asphériques? Si c'était faux, je songe à un axiome de `descente homotopique', qui impliquerait qu'un objet $W_\infty$-asphérique est aussi $W$-asphérique.

% foot: Version 7 décembre 2002
\printedpage{11}
% head: A. Grothendieck : Dérivateurs, Ch. XII, Édition des Universités de Montpellier II et Paris VII

Je note dès maintenant que si des $W_i$ ($i \in I$) satisfont aux axiomes W(1) à W(5), il en est de même de $\bigcap_i W_i$ (\footnote[30]{C'est vrai pour tout ensemble d'axiomes dans la nature d'axiomes de stabilité.}). Donc il y a un plus petit ensemble de flèches $W_\omega$ satisfaisant à ces conditions. S'il n'est égal à $W_\infty$, i.e. si les axiomes W(1) à W(5) ne caractérisent déjà $W_\infty$, ce serait un localiseur vraiment très intéressant.

Je me rends compte que je n'ai utilisé qu'une forme assez faible de l'axiome W(4) - essentiellement, la validité de la proposition 4. Je vais résumer différentes variantes de l'axiome de fibration.

\textbf{W(4)}

\begin{itemize}[leftmargin=4em, labelwidth=2.5em, labelsep=0.6em, align=right]
\item[a$_1$)] Si $f : X \longrightarrow Y$ dans Cat est une $W$-fibration gauche, c'est une $W$-fibration ultraforte.

\item[a$_2$)] Si $f : X \longrightarrow Y$ est une $W$-fibration droite, c'est une $W$-fibration ultraforte.

\item[a$'$)] Si $f : X \longrightarrow Y$ est $W$-parfait, i.e. à la fois $W$-propre et $W$-lisse (\footnote[31]{ce qui implique que c'est une $W$-fibration gauche \textit{et} droite.}), alors $f$ est une $W$-fibration ultraforte.

\item[a$''$)] Si $f : X \longrightarrow Y$ est à la fois homotopiquement propre et homotopiquement lisse (cf. NB page 10), $f$ est une $W$-fibration ultraforte. (Il vaut mieux dire : Si $f$ est $W_\omega$-parfait, c'est une $W$-fibration ultraforte.)

\item[b)] Soit $f : X \longrightarrow Y$ une flèche dans Cat$/S$, avec $X$ et $Y$ des $W$-fibrés sur $S$. Si $f \in W$, alors $\forall\, s \in S$, $f_s : X_s \longrightarrow Y_s$ est \texttt{[dans]} $W$.

\item[b$'$)] Soit $f : X \longrightarrow Y$ une flèche de Cat$/S$ qui soit $W$-parfaite et $W_\infty$-parfaite, alors $f \in W$ implique que les fibres de $f$ sont $W$-asphériques. (Il vaut mieux dire : Si $f$ est $W_\omega$-parfait, il est une $W$-fibration ultraforte.) \texttt{[Plutôt : Soit $f : X \longrightarrow Y$ une flèche dans \textrm{Cat}$/S$ telle que $X \longrightarrow S$ et $Y \longrightarrow S$ soient $W_\omega$-parfaites. Si $f \in W$, alors $\forall\, s \in S$, $f_s : X_s \longrightarrow Y_s$ est dans $W$.]}

\item[b$''$)] \msp{page 14}Soit $W_\omega \subseteq W \cap W_\infty$ la plus petite partie de Fl(Cat) satisfaisant les conditions W(1,2,3) (\footnote[32]{Il vaut mieux prendre W(2 bis) au lieu de W(2), pour trouver $W_\omega$ la plus petite possible!}) (axiomes de saturation, de localisation, de l'objet final). Soit $f : X \longrightarrow Y$ une flèche dans Cat qui soit $W_\omega$-parfaite. Si $f \in W$, alors les fibres de $f$ sont $W_\omega$-asphériques \texttt{[plutôt $W$-asphériques]}.
\end{itemize}

Entre ces propriétés de $W$, on a les implications tautologiques (sauf a$''$) $\Longrightarrow$ b), qui est essentiellement la proposition 4)

% foot: Version 7 décembre 2002
\printedpage{12}
% head: A. Grothendieck : Dérivateurs, Ch. XII, Édition des Universités de Montpellier II et Paris VII

\[
\begin{tikzcd}[row sep=huge, column sep=0pt]
\mathrm{a} \arrow[d, Rightarrow] & {}= (\mathrm{a}_1 \Longleftrightarrow \mathrm{a}_2) & \hspace{2.5cm} \\
\mathrm{a}' \arrow[d, Rightarrow] & & \smash{\left.\rule[-1.8cm]{0pt}{3.6cm}\right\}\ \text{\begin{tabular}{l}En fait,\\ toutes ces\\ conditions sont\\ équivalentes\end{tabular}}} \\
\mathrm{a}'' \arrow[d, Rightarrow] & & \\
\mathrm{b} \arrow[d, Rightarrow] & & \\
\mathrm{b}' \arrow[d, Rightarrow] & & \\
\mathrm{b}'' & &
\end{tikzcd}
\]

(\footnote[33]{Je reviens sur les diverses variantes de W(4) plus bas, p. 53-56.}).

Eh bien, dans le lemme page 12, il suffisait de supposer que $W$ satisfait à la version la plus faible b$''$) de l'axiome des fibrations (sept versions!). En effet, les fibrations $\bar{X} \longrightarrow Y$ et $E \longrightarrow X$ utilisées dans la démonstration du lemme sont $W_\omega$-parfaites, et il suffisait de savoir que si elles sont dans $W$, alors les fibres sont $W$-asphériques.

Mais je ne pense pas que même la version la plus forte W(4a) de l'axiome des fibrations, jointe aux trois axiomes précédents et à W(5) (connexité), implique $W = W_\infty$. Il faudrait tout au moins y joindre l'axiome

\begin{itemize}[leftmargin=4em, labelwidth=3em, labelsep=0.6em]
\item[\textbf{W(6)}] (Axiome des limites filtrantes.) Soit $X = \varinjlim\limits_{i} X_i$ une limite inductive filtrante dans Cat. Si les $X_i$ sont $W$-asphériques, $X$ aussi.
\end{itemize}

Ceci devrait être équivalent à la propriété, en apparence plus forte, en termes de flèches de $W$ :

\textbf{Proposition 5.} \textit{Supposons qu'on ait} W(6) (en plus de W(1) à W(4) (\footnote[34]{W(4) sous la forme la plus faible W(4b$''$) .})),
\msp{page 15}\textit{et soit}
\[ (f_i)_{i\in I} : (X_i)_{i\in I} \longrightarrow (Y_i)_{i\in I} \]
% foot: Version 7 décembre 2002
\printedpage{13}
% head: A. Grothendieck : Dérivateurs, Ch. XII, Édition des Universités de Montpellier II et Paris VII
\textit{un homomorphisme de systèmes inductifs dans} Cat\textit{, indexé par un ensemble ordonné filtrant, soient $X = \varinjlim X_i$, $Y = \varinjlim Y_i$, $f = \varinjlim f_i$, $f : X \longrightarrow Y$. Si $f_i$ est dans $W$ pour tout $i$, alors $f \in W$.}

On factorise chaque $f_i$ en
\[ X_i \xrightarrow{\alpha_i} \bar{X}_i \xrightarrow{\bar{f}_i} Y_i, \qquad \alpha_i \in W_\omega,\ \bar{f}_i \in \mathrm{Parf}\, W_\omega \]
(\footnote[35]{$W_\omega$ introduit dans la condition W(4b$''$), page 14.}) suivant le système breveté. La construction étant fonctorielle, les $\bar{X}_i$ forment eux aussi un système inductif. De plus, comme le foncteur $S \longrightarrow \Ch(S)$ qui donne naissance à la factorisation canonique des morphismes, commute aux $\varinjlim$ filtrantes (immédiat), il s'ensuit que, si on pose
\[ \bar{X} = \varinjlim_{i} \bar{X}_i \]
et qu'on regarde les flèches
\[ X \xrightarrow{\alpha} \bar{X} \xrightarrow{\bar{f}} Y, \qquad \alpha = \varinjlim \alpha_i,\ \bar{f} = \varinjlim \bar{f}_i, \]
avec $\bar{f}\alpha = \bar{f}$, cette factorisation de $f$ est bien sa factorisation canonique. Donc $\alpha \in W$, % [sic] $\bar{f}\alpha = \bar{f}$ -> $\bar{f}\alpha = f$, the composite should be f
$\bar{f} \in \mathrm{Parf}\, W_\omega$. Ceci vu, en vertu de prop. 3, $f \in W$, i.e. $\bar{f} \in W$, est vrai si les fibres de $\bar{f}$ sont $W$-asphériques. Mais les fibres de $\bar{f}$ s'interprètent comme limites inductives filtrantes de fibres des $\bar{f}_i$, lesquelles sont $W$-asphériques par la condition W(4b$''$) des fibrations. Donc l'axiome W(6) donne ce qu'il faut.

\textbf{Conjecture} (hardie? sotte?). \textit{Si $W \subseteq$} Fl(Cat) \textit{satisfait les axiomes} W(1) \textit{à} W(6) \textit{(avec} W(4) \textit{sous sa forme la plus faible} W(4b$''$) (\footnote[36]{et peut-être W(2) sous la forme plus faible W(2 bis).})\textit{), alors $W = W_\infty$.}

\msp{page 16}Mais il est bien possible qu'il faille renforcer un peu W(1) (saturation), en exigeant que tout rétracte d'un $f \in W$ soit dans $W$.

\textbf{Remarque.} Considérons, pour $n \in \mathbf{N}$, l'ensemble $W_n \subseteq$ Fl(Cat) des `$n$-équivalences d'homotopie' dans Cat :
\[ f \in W_n \iff \pi_0(f) \text{ bijectif, et } \pi_i(f, x) : \pi_i(X, x) \longrightarrow \pi_i(Y, f(x)) \text{ bijectif pour } 0 < i \leq n. \]

Je croyais que ça satisfaisait à tout ce qu'on peut demander à un localiseur, sans avoir pris la peine de l'écrire avec soin. Ma ça ne peut satisfaire à W(1) jusqu'à W(4b$''$), puisqu'on % [sic] Ma -> Mais, word truncated
n'a pas $W_n \subseteq W_\infty$!

W(1) est trivial sur la définition (et on a sans larmes même la \textit{saturation forte} de $W_i$, % [sic] $W_i$ -> $W_n$ (also in the following lines), the remark concerns W_n
i.e. si $f$ dans Cat est \texttt{[un]} isomorphisme dans $\mathrm{Hot}_{W_i} = (\mathrm{Cat})W_i^{-1}$, alors $f \in W_i$). Je crois me rappeler avoir vérifié W(2), par la même méthode que pour $W_\infty$, en utilisant un critère cohomologique à la \textsc{Artin-Mazur}, pour exprimer $f \in W_i$, et en utilisant les suites spectrales de Leray pour $X \longrightarrow S$ et $Y \longrightarrow S$. La condition W(3) est trivialement satisfaite puisque $W_n \supseteq W_\infty$ et qu'elle l'est pour $W_\infty$. Donc ça canule là où je m'y
% foot: Version 7 décembre 2002

\end{document}
